(# The generated parts of the DRLD in one document, rendered by generate_drld.py with `doc` set to a Document:
   the data-item paragraphs (DRLD chapter "DRL Data Items and Structures"), the recipe cards ("Pipeline Recipes")
   and the QC-parameter cards ("QC Parameters"), each grouped as the DRLD groups them. With `doc.standalone` the
   file compiles on its own, provided the DRLD sources are reachable through TEXINPUTS (it reuses the DRLD's
   normal_style, styles_data and acronyms); otherwise it is a fragment to \input into the DRLD. #)
(% if doc.standalone %)
\documentclass[11pt,oneside,a4paper]{article}
\raggedbottom
\textheight=246mm
\textwidth=164mm
\oddsidemargin=-2mm
\parindent=0mm
\parskip=0.3em
\usepackage[T1]{fontenc}
\usepackage{mathptmx}
\usepackage[scaled=0.92]{helvet}
\usepackage[dvipsnames]{xcolor}
\usepackage{colortbl}
\usepackage{longtable}
\usepackage{caption}
\usepackage{array}
\usepackage{tabularx}
\usepackage{booktabs}
\usepackage{enumitem}
\usepackage{listings}
\usepackage{tikz}
\usepackage[many]{tcolorbox}
\usepackage{xstring}
\usepackage{hyperref}
\usepackage[nohyperlinks]{acronym}
\setcounter{secnumdepth}{4}
\setcounter{tocdepth}{2}
\lstset{basicstyle=\ttfamily, columns=flexible, showspaces=false}
\input{normal_style}
\input{styles_data}
\providecommand{\CODE}[1]{\lstinline[]!#1!}
\providecommand{\DRL}[1]{\lstinline[style=DRLstyle]!#1!}
\providecommand{\TPL}[1]{\lstinline[style=TPLstyle]!#1!}
\providecommand{\REQ}[1]{\texttt{#1}}
\newenvironment{recipedef}
  {\begin{tcolorbox}[breakable,enhanced jigsaw]\begin{longtable*}{p{0.24\hsize}p{0.69\hsize}}}
  {\end{longtable*}\end{tcolorbox}}
\newenvironment{datastructdef}
  {\begin{tcolorbox}[breakable,enhanced jigsaw]\begin{longtable*}{p{0.97\hsize}}}
  {\end{longtable*}\end{tcolorbox}}
\title{METIS DRLD --- generated cards}
\author{generate\_drld.py, pymetis (* doc.version|latex *)}
\date{(* doc.date *)}
\begin{document}
\maketitle
\tableofcontents
\clearpage
\input{acronyms}
\clearpage
(% endif %)
(# ---------------------------------------------------------------- data items #)
\section{DRL Data Items and Structures}\label{sec:drl_data_structures}
(% for family in doc.item_families %)
\subsection{(* family.title|latex *) data items}
(% for kind in family.kinds %)
\subsubsection{(* kind.title|latex *)}
(% for item in kind.items %)
(% include 'dataitem.tex' %)
(% endfor %)
(% endfor %)
(% endfor %)
(# ---------------------------------------------------------------- recipes #)
\clearpage
\section{Pipeline Recipes}\label{sec:pipeline_recipes}
(% for family in doc.recipe_families %)
\subsection{(* family.title|latex *)}
(% for recipe in family.recipes %)
\subsubsection{\REC*{(* recipe.name *)}: (* recipe.synopsis|latex *)}\label{rec:(* recipe.name *)}
(* recipe.description|latex *)

(% include 'recipe.tex' %)
(% endfor %)
(% endfor %)
(# ---------------------------------------------------------------- QC parameters #)
\clearpage
\section{QC Parameters}\label{sec:qc_parameters}
Placeholders in a parameter name (\CODE{band}, \CODE{cgrph}, \CODE{nn}) stand for values known only at run time.
(% for family in doc.qc_families %)
\subsection{(* family.title|latex *)}
(% for qc in family.qcs %)
(% include 'qc.tex' %)
(% endfor %)
(% endfor %)
(% if doc.standalone %)
\end{document}
(% endif %)
